% =============================================================================
% Annexes
% =============================================================================
\chapter{Schéma de sortie d'une section}
\label{annexe:schema}

Le schéma ci-dessous est le contrat entre le générateur et le reste du
système. Toute sortie qui s'en écarte est rejetée par la validation
déterministe, avant même la revue humaine.

\begin{lstlisting}[style=btpjson, caption={Contrat JSON d'une section générée}]
{
  "section_key": "memoire_technique",
  "titre": "Memoire technique d'execution",
  "blocs": [
    {
      "type": "paragraphe",
      "texte": "La methode d'execution retenue s'appuie sur..."
    },
    {
      "type": "liste",
      "texte": "- Phase preparatoire\n- Phase chantier\n- Reception"
    }
  ],
  "citations": ["chunk-001", "chunk-014", "chunk-027"],
  "exigences_couvertes": ["ex-012", "ex-018"],
  "abstentions": [
    "Aucune reference patrimoniale superieure a 20 M EUR n'a ete trouvee."
  ]
}
\end{lstlisting}

Les champs \texttt{citations} et \texttt{exigences\_couvertes} sont des
identifiants, jamais des libellés libres. C'est cette contrainte qui rend
la citation auditable.

\chapter{Protocole du jeu de référence}
\label{annexe:golden}

Le jeu de référence utilisé pour l'évaluation de la récupération
comprend trente questions, réparties en trois familles.

\begin{btptable}{Composition du jeu de référence}
\begin{tabularx}{\linewidth}{@{}L{4.2cm} C{2.2cm} Y@{}}
\toprule
\rowcolor{btppale}
\btpth{Famille} & \btpth{Nombre} & \btpth{Exemple} \\
\midrule
Terme exact / identifiant & 10 & « Que dit l'article 12.3 du CCAP~? » \\
Concept métier paraphrasé & 10 & « Quelles sanctions en cas de retard~? » \\
Question mixte            & 10 & « Quelles références en assainissement, norme ISO 14001~? » \\
\bottomrule
\end{tabularx}
\end{btptable}

Pour chaque question, deux à cinq extraits pertinents ont été identifiés
manuellement par un praticien. Un extrait est jugé pertinent s'il
contient une information nécessaire à une réponse complète, pas seulement
s'il partage des mots avec la question.

\chapter{Glossaire des pièces d'un dossier}
\label{annexe:pieces}

Ce glossaire complète le tableau~\ref{tab:pieces-dao}. Il n'a pas vocation
à se substituer aux définitions réglementaires.

\begin{description}[style=nextline, leftmargin=1.2em, labelindent=0pt, font=\bfseries\color{btpdeep}]
  \item[Règlement de consultation.] Pièce qui fixe les règles du jeu~:
        calendrier, modalités de remise, critères de jugement, pièces
        exigées.
  \item[CCAG.] Clauses administratives générales applicables à une
        famille de marchés (travaux, fournitures, prestations
        intellectuelles).
  \item[CCAP.] Clauses administratives particulières au marché, qui
        précisent ou dérogent au CCAG.
  \item[CCTP.] Spécifications techniques de l'ouvrage ou de la
        prestation.
  \item[DPGF.] Décomposition du prix global et forfaitaire, tableau
        structuré des coûts.
  \item[BPU.] Bordereau des prix unitaires, utilisé notamment dans les
        marchés à prix unitaires.
\end{description}

\btpEndMark
\clearpage
